\begin{remark}[7.1.6]
\label{IV.2.s7.1.6-fr}
Se ha visto (6.3.8) que, bajo las hipótesis de
(7.1.5), para todo anillo cociente $B$ de $A$,
las fibras del morfismo canónico
$\operatorname{Spec}(\widehat B)\to\operatorname{Spec}(B)$
son esquemas de Cohen-Macaulay. Por consiguiente,
según (6.4.3) y (5.7.5), \emph{para que $B$ no tenga
ideales primos asociados inmersos es necesario y
suficiente que $\widehat B$ tampoco los tenga}.
\end{remark}
